% ============================================================
% Data-quality report: NYC TLC High-Volume FHV trip records, 2020-04 to 2026-07
% Compiles from the repo root:  latexmk -pdf docs/data_quality/hvfhv_data_quality.tex
% Tables/figures/numbers are generated by code/build/make_dq_report_tables.py
% (output/sum/dq_*, output/fig/dq_*, output/sum/dq_numbers.tex) -- do not edit by hand.
% ============================================================
\documentclass[11pt]{article}
\usepackage[margin=1in]{geometry}
\usepackage{graphicx,adjustbox,booktabs,amsmath,setspace,caption,array}
\usepackage{natbib}
\usepackage[colorlinks=true,allcolors=blue]{hyperref}
\graphicspath{{output/fig/}}
\setcitestyle{authoryear,round}
\emergencystretch=3em
\input{output/sum/dq_numbers}
\newcommand{\dqtable}[1]{\begin{adjustbox}{max width=\linewidth,center}\input{output/sum/#1}\end{adjustbox}}

\title{Data Quality of the NYC TLC High-Volume For-Hire Vehicle Trip Records, April 2020--July 2026}
\author{Prepared by EK}
\date{\today}

\begin{document}
\maketitle

\section{Summary}
This report audits the full TLC High-Volume FHV (HVFHV) trip-record series from April 2020 through July 2026
(\dqNmonths{} monthly files, \dqTotalTrips{} trip records). Everything is computed on the raw files with no
cleaning filters unless a panel says otherwise. All timestamps are local New York wall-clock time.

\paragraph{1. Is any data missing?}
\begin{itemize}
  \item \textbf{Months, days, hours.} All \dqNmonths{} monthly files are present with no gap in the month
  sequence. Section~\ref{sec:complete} lists the \dqNFlaggedCity{} citywide days that fall below 50\% of their
  centered 28-day median or have zero trips; \dqNFlaggedUnexplained{} of them are not explained by a documented event
  or holiday. There are \dqNZeroHours{} hours with zero citywide trips.
  \item \textbf{Platforms.} Uber and Lyft appear throughout. Via's records run from \dqFirstVia{} to \dqLastVia{}, ending abruptly after \dqViaLastDay{} (about \dqViaAvgDaily{} trips per day before then), although Via stated that it ran its New York City rideshare service until 2021-12-20 \citep{via_letter_2021}, so roughly ten weeks of Via operation are not in the file; the cause is unknown. Juno has
  \dqTripsJuno{} records in this period; unrecognized license codes (``other'') account for \dqTripsOther{} records.
  \item \textbf{Fields.} Two nulls are structural: \texttt{cbd\_congestion\_fee} is \dqPctNullCbd{} null over the full
  period because it first appears in 2025, and \texttt{airport\_fee} is \dqPctNullAirportFee{} null
  (a null-typed column in 2020-04 and 2020-10, and null through part of 2021). Three fields are missing for
  \emph{reporting-practice} reasons that differ by platform and period (Table~\ref{tab:missing_platform}):
  \texttt{on\_scene\_datetime} is populated for Uber but null for \dqNullOnSceneLyft{} of Lyft records overall
  (\dqOnSceneLyftBefore{} before 2025) and \dqNullOnSceneVia{} of Via records, so it is not
  ``Accessible Vehicles only'' as the dictionary states; \texttt{originating\_base\_num} is blank for
  \dqNullOrigBaseLyft{} of Lyft and \dqNullOrigBaseVia{} of Via records; and \texttt{access\_a\_ride\_flag} is blank
  (not ``N'') for \dqNullAarBlankUber{} of Uber records, through \dqAarBlankLastMonth. Every other numeric field and
  timestamp has a null share of at most \dqMaxIncidentalNull{}.
\end{itemize}

\paragraph{2. Does every ride from these companies end up in the data?}
Not necessarily, and the reasons are partly by design (Section~\ref{sec:legal}). The records cover trips
dispatched through HVFHS bases; cancelled requests, trips dispatched by other bases, and taxi trips booked
through ride-hail apps are not in this file. Against TLC's own published aggregates, the record counts differ by
\dqBaseGapMin{} to \dqBaseGapMax{} per year against the base aggregate report and by \dqIndGapMin{} to \dqIndGapMax{}
against the industry indicators; \dqNBaseFlagged{} of \dqNBaseMonths{} platform-months differ by more than 2\%
(largest: \dqBaseWorstGap{} in \dqBaseWorstMonth), and \dqIndFlagged{} of \dqIndMonths{} months differ by more than 2\%
against the industry indicators (\dqIndWorstGap{} in \dqIndWorstMonth). Both TLC series are derived from the same base
submissions, so this is a consistency check, not an independent count of all rides: it shows the public file reproduces
what bases reported, not that bases reported everything.

\paragraph{3. Summary statistics.} Table~\ref{tab:summary_full} gives mean, standard deviation, median, minimum,
maximum and \% missing for every main variable over the whole period, raw (Panel A) and after the cleaning filters (Panel B).
The filters keep \dqPctPassFilters{} of records overall.

\section{Data and coverage}
The source is the TLC monthly HVFHV parquet files (\texttt{fhvhv\_tripdata\_YYYY-MM.parquet}), downloaded
2026-09-26 and checked against the download manifest (byte sizes match). The series starts in 2019-02; this report covers
2020-04 onward by choice. The sum of per-month record counts in the audit equals the sum of the parquet footer row counts,
so no row was silently dropped by the audit. Schema drift handled in the audit: \texttt{airport\_fee} is a null-typed
column in 2020-04 and 2020-10; from 2023-02 strings are stored as large strings and location IDs as 32-bit integers;
from 2025-01 a 25th column, \texttt{cbd\_congestion\_fee}, appears. Field definitions follow the TLC data dictionary
\citep{tlc_dictionary}. TLC states, for its taxi data, that the trip data were not created by the TLC and that it cannot guarantee their
accuracy \citep{tlc_userguide}, and for FHV base data that it publishes them as submitted by the bases and cannot guarantee their accuracy or completeness \citep{tlc_base_aggregate}; this is the reason the audit reports anomalies instead of assuming clean inputs.

\section{What the law requires reported, and what is excluded}\label{sec:legal}
\paragraph{Who reports.} Local Law 149 of 2018 defines a high-volume for-hire service as a business, operating under
one brand or trade name, that connects passengers to for-hire vehicles through software and dispatches 10,000 or more
trips in the city in one day; all bases under a common brand are counted together \citep{ll149_2018}. TLC's user guide
records that the HVFHS license class was created in 2019 and that, as of February 2019, the qualifying companies were
Uber, Lyft, Via and Juno \citep{tlc_userguide}.

\paragraph{What must be reported.} TLC's rules in Chapter~59 of Title~35 of the Rules of the City of New York,
section 59D-14, require an HVFHS to transmit, for the trips it dispatches, the pickup and drop-off date, time and
location; the driver and vehicle license numbers; the affiliated base; passenger count; trip mileage; the date and time the
passenger requested the trip; the itemized fare (fare, tolls, surcharges, commission rate, deductions, gratuity); the
payment the driver received; and an indicator for MTA Access-A-Ride trips \citep{tlc_rules_2021}. The 2021 amendment
also added driver availability records, weekly driver earnings statements, and trip-level time and mileage outside
city limits \citep{tlc_rules_2021}. The public HVFHV files expose a subset of these fields and omit driver identifiers.

\paragraph{What is excluded by design.}
\begin{itemize}
  \item \emph{Trips dispatched by other bases.} The FHV data include high-volume bases and also community livery, luxury
  limousine and black-car bases \citep{tlc_userguide}; the HVFHV file contains only the high-volume license class, so trips
  dispatched through smaller bases are not in it.
  \item \emph{Taxi trips booked through apps.} Yellow taxis can be hailed through e-hail apps such as Curb and Arro
  \citep{tlc_userguide}, and Uber announced in 2022 that New York taxis would be integrated into its app \citep{uber_taxi_thedrive}. Taxi trips are
  recorded in the taxi files, so such trips would not be in the HVFHV file. We have not verified how TLC classifies them.
  \item \emph{Cancelled or unfulfilled requests.} The rule text refers to trips the service dispatches \citep{tlc_rules_2021}, and TLC's user guide says it
  collects records for each trip \emph{completed} by licensed drivers and vehicles \citep{tlc_userguide}. We found no
  explicit statement on how cancelled requests are treated, and the files contain no cancellation field; the
  documentation points to completed or dispatched trips only, but this is not stated outright by a primary source.\textsuperscript{?}
  \item \emph{Privacy.} The public files report pickup and drop-off at the taxi-zone level and omit driver identity.
  \item \emph{Access-A-Ride.} The rules require an indicator for MTA Access-A-Ride trips \citep{tlc_rules_2021}, and
  these trips are included and flagged (\texttt{access\_a\_ride\_}\allowbreak\texttt{flag});\dqNAarY{} records carry ``Y'' (\dqPctAarY{} of
  all records). The flag is blank for \dqNAarBlank{} Uber records through \dqAarBlankLastMonth, so Access-A-Ride trips
  cannot be identified for Uber before then.
\end{itemize}
TLC's base aggregate report carries its own caveat: TLC publishes base trip record data as submitted and cannot
guarantee its accuracy or completeness \citep{tlc_base_aggregate}.

\paragraph{Platform entry and exit.} Juno stopped operating in New York in November 2019 \citep{juno_techcrunch_2019},
before this sample, consistent with \dqTripsJuno{} Juno records here. Via's first and last month in the records are
\dqFirstVia{} and \dqLastVia{} (Table~\ref{tab:platform_year}); its last day with a full volume of trips is \dqViaLastDay{}.
Via announced that it would discontinue its New York City rideshare service on 2021-12-20 \citep{via_letter_2021}, so records for
the final weeks of Via's operation are absent from the HVFHV file. We cannot tell from the data whether Via stopped
reporting, moved those trips to another license class, or reported them elsewhere; this is an open question for TLC.

\section{Completeness}\label{sec:complete}
\begin{table}[ht]\centering\small
\caption{Trip records by platform and year}\label{tab:platform_year}
\dqtable{dq_trips_platform_year}
\end{table}

\begin{figure}[ht]\centering
\includegraphics[width=\linewidth]{dq_monthly_trips_platform.png}
\caption{Monthly trip records by platform, with the sum of TLC's base aggregate report for the same bases (dashed).}
\label{fig:monthly}
\end{figure}

\paragraph{Benchmark against TLC aggregates.} Platform bases were mapped to companies using the
\texttt{dispatching\_base\_num} values observed in the trip data. Table~\ref{tab:benchmark} compares record counts
with the FHV Base Aggregate Report \citep{tlc_base_aggregate} and with the HVFHS trips per day in TLC Industry Indicators
\citep{tlc_industry_indicators} (trips per day times days in month; available through 2026-05). Monthly detail is in
\texttt{output/sum/dq\_benchmark\_monthly.csv}.
\begin{table}[ht]\centering\small
\caption{Records versus TLC aggregates, by year}\label{tab:benchmark}
\dqtable{dq_benchmark}
\end{table}

\paragraph{Missing days and hours.}
Figure~\ref{fig:daily} marks citywide days that are zero or below 50\% of the centered 28-day median.
Table~\ref{tab:flagged} lists them with an explanation where a documented event or holiday applies. A dip is labelled
``explained'' only if it matches a listed event; everything else is ``unexplained'' and should not be read as missing data
or as a real event without further checking. Events used for labelling are listed in \texttt{make\_dq\_report\_tables.py}; the 2026 winter storms are documented in \citet{fern_nbn_2026,blizzard_aljazeera_2026}. The \dqNZeroHours{} hours with zero citywide trips all fall between
\dqZeroHoursFirst{} and \dqZeroHoursLast{}, in the June 2020 curfew window; \dqZeroHoursExplained{} of them match that
event. Platform-level flags (including Via's zero-trip days after its last month in the data) are in
\texttt{output/sum/dq\_flagged\_days\_all.csv}.
\begin{figure}[ht]\centering
\includegraphics[width=\linewidth]{dq_daily_trips_flagged.png}
\caption{Citywide trip records per day; flagged days marked.}\label{fig:daily}
\end{figure}
\begin{table}[ht]\centering\scriptsize
\caption{Flagged citywide days}\label{tab:flagged}
\dqtable{dq_flagged_days}
\end{table}

\section{Missing fields}
\begin{figure}[ht]\centering
\includegraphics[width=\linewidth]{dq_missing_heatmap.png}
\caption{Percent of records missing, by numeric field and month.}\label{fig:heatmap}
\end{figure}
\begin{table}[ht]\centering\small
\caption{Percent missing by field and year}\label{tab:missing}
\dqtable{dq_missing_by_year}
\par\footnotesize $^{\dagger}$ Structural, not incidental: \texttt{cbd\_congestion\_fee} first appears in 2025-01;
\texttt{airport\_fee} is a null-typed column in the 2020-04 and 2020-10 files. $^{\ddagger}$ Missingness is specific to a
platform and period (Table~\ref{tab:missing_platform}). Blank strings in categorical fields are counted as missing.
\end{table}

\begin{table}[ht]\centering\small
\caption{Percent missing for platform-specific fields, by platform and year}\label{tab:missing_platform}
\dqtable{dq_missing_platform}
\par\footnotesize ``--'' means the platform has no records that year. Blank strings count as missing.
\end{table}

\section{Summary statistics}
\begin{table}[ht]\centering\scriptsize
\caption{Summary statistics over the whole period, 2020-04 to 2026-07}\label{tab:summary_full}
\dqtable{dq_summary_full}
\par\footnotesize Panel A: all raw records. Panel B: after dropping fare $\le 0$ or $>\$1{,}000$; miles $\le 0$ or $>100$;
trip time $\le 60$ s or $>6$ h; unknown zones (264, 265) or null zone; Access-A-Ride trips. Mean, SD, min, max and
\% missing are exact. Median, P1 and P99 are exact pooled quantiles (DuckDB \texttt{quantile\_cont}) over all 76 files, one query per variable. \texttt{wait\_time} = pickup $-$ request (seconds);
\texttt{driver\_share} = \texttt{driver\_pay}/\texttt{base\_passenger\_fare}; \texttt{platform\_margin} = fare $-$ driver pay.
\end{table}

\begin{table}[ht]\centering\scriptsize
\caption{Means and medians by year, headline variables}\label{tab:by_year}
\dqtable{dq_summary_by_year}
\par\footnotesize Yearly medians are trip-weighted averages of monthly medians (approximate); means are exact.
\end{table}

\begin{figure}[ht]\centering
\includegraphics[width=\linewidth]{dq_fare_share_medians.png}
\caption{Monthly median fare per mile and driver share, by platform (raw records).}\label{fig:fare_share}
\end{figure}

\section{Anomalies and recommended cleaning rules}
Table~\ref{tab:oor} reports each rule as a percent of records by year. Overall, \dqNFareLeZero{} records have a fare
at or below zero (\dqPctFareLeZero), \dqNMilesLeZero{} have zero miles, \dqNTimeLeSixty{} have trip time of 60 seconds or
less, and \dqNDoUnknown{} (\dqPctDoUnknown) have an unknown drop-off zone (264/265). There are \dqDupTotal{} exact duplicate
rows. \dqNRequestAfterPickup{} records (\dqPctRequestAfterPickup) have a request time after the pickup time, which gives a
negative \texttt{wait\_time}. \dqPctDriverShareGtOne{} of records have \texttt{driver\_share} above 1 (driver pay exceeding the
fare); by year this ranges from \dqDsGtOneMin{} (\dqDsGtOneMinYear) to \dqDsGtOneMax{} (\dqDsGtOneMaxYear). Because
\texttt{driver\_pay} is net of commission and excludes tips, a share above 1 means the platform paid more than the base
fare, and it is not a unit error we can fix by filtering; we report it rather than drop it.

\paragraph{Do the cleaning rules bite evenly over time?} The share of records passing the filters ranges from
\dqPassMin{} (\dqPassMinYear) to \dqPassMax{} (\dqPassMaxYear). A filter that removes much more in some periods
changes the sample composition, so apply the same rules in every period and report drops by window.

\begin{table}[ht]\centering\scriptsize
\caption{Anomaly rules, percent of records by year}\label{tab:oor}
\dqtable{dq_out_of_range}
\end{table}

\paragraph{Recommended rules.} Keep the CLAUDE.md project filters: fare in $(0,1000]$; miles in $(0,100]$; trip time
in $(60\text{ s}, 6\text{ h}]$; drop zones 264/265; consider excluding Access-A-Ride. Set negative \texttt{wait\_time} to missing
rather than dropping the trip. Do not filter on \texttt{on\_scene\_datetime}, and do not build a wait-time measure from it: it is populated for one platform
only for most of the sample. Note that an Access-A-Ride exclusion based on \texttt{access\_a\_ride\_flag} removes flagged
trips for Lyft and Via throughout but for Uber only from \dqAarBlankLastMonth{} onward, so it would filter
asymmetrically across platforms and time; report results with and without it.

\appendix
\section{Post-flight verification (CoVe)}
\input{docs/data_quality/cove_block}

\bibliographystyle{apalike}
\bibliography{docs/data_quality/references}
\end{document}
